\documentclass[leqno]{jsarticle}
\usepackage{amssymb}
\usepackage{amsthm}
\usepackage{amsmath}
\usepackage{comment}
	%可換図式用
		%\usepackage[all]{xy}
	%重くなるので使わいときはコメント化しておく。
%定理環境 thmはあまり使わずpropをなるべく使う。
%主張は基本的に証明の中で使い、そのため番号を付けない。
%注意環境を付けてみた。
\newtheorem{thm}{定理}
\newtheorem{prop}[thm]{命題}
\newtheorem{cor}[thm]{系}
\newtheorem{lem}[thm]{補題}
\newtheorem{df}[thm]{定義}
\newtheorem{exam}[thm]{例}
\newtheorem{fact}[thm]{事実}
\newtheorem*{claim}{主張}
\newtheorem*{cau}{注意}
\renewcommand{\proofname}{{\bf 証明}}
%矢印たち。
%\toは\rightarrowと同じ。\getsは\leftarrowと同じ。同値は\iff
%上向き矢はuparrow逆はdownarrow
\newcommand{\To}{\Longrightarrow}
\newcommand{\Gets}{\Longleftarrow}
%黒板太字
\newcommand{\nn}{\mathbb{N}}
\newcommand{\zz}{\mathbb{Z}}
\newcommand{\qq}{\mathbb{Q}}
\newcommand{\rr}{\mathbb{R}}
\newcommand{\cc}{\mathbb{C}}
%無限大
\newcommand{\mgn}{\infty}
%ギリシャン
\newcommand{\dpst}{\displaystyle}
\newcommand{\ep}{\varepsilon}
\newcommand{\al}{\alpha}
\newcommand{\be}{\beta}
\newcommand{\de}{\delta}
\newcommand{\la}{\lambda}
\newcommand{\La}{\Lambda}
\newcommand{\ga}{\gamma}
\newcommand{\lil}{\la\in \La}
%商集合
\newcommand{\qt}[2]{#1/\! #2}
%enumerate環境
\renewcommand{\labelenumi}{{\rm (\arabic{enumi})}}
%以降alignじゃなくてenumerate を使え。
%なんか演算
\newcommand{\card}{\text{{\rm card}}}
\newcommand{\supp}{\text{{\rm supp}}}
%なんか演算２
\newcommand{\bd}{\text{{\rm Bdry}}}
\newcommand{\cl}{\text{{\rm CL}}}
\newcommand{\nai}{\text{{\rm Int}}}
\newcommand{\ind}{\text{{\rm ind}}}
%差集合は\setminusで統一する。等号の否定は\neq
%真部分集合は\subsetneqq　偏微分は\partial
%ディスプレイスタイル。\dfracでディスプレイ分数
%空集合は\Oを使う！！！！！！！！！！！！

\title{次元論のひとつの「補題」}
\author{一色三良(concious77)}
\date{\today}
			\begin{document}
\maketitle
この文章では参考文献\cite{0}に紹介されている多様体の埋め込み定理のための次元論的な補題、この文章で定理\ref{main}とされている命題を証明することである。しかし、次元論を最初から紹介するのはとても面倒なことなので、詳しくは参考文献\cite{-1}、\cite{1}を当たってほしい。以下の文章用いられる次元論的な命題は事実\ref{00}、\ref{fact}、\ref{qq}に限られることを注意しておく。つまり、これだけの事実が分かっていればこの文章を読むには十分であるということである。\par
以下の文章では\textbf{すべての空間を可分距離空間とする。}よって、この文章で単に「空間」と言ったら可分距離空間のことを指す。\par
\textbf{第二可算な位相多様体は可分な距離付け可能空間である}という事実を知っておけば定理\ref{main}が第二可算多様体に適用することができ、参考文献\cite{0}の多様体の埋め込み定理の大きな行間を埋めることができることであろう。\par
多様体は局所コンパクトハウスドルフ空間であり、特に正則空間なので、第二可算な多様体は\textbf{ウリゾーンの距離付け可能定理から距離付け可能である。}ことが分かる。また第二加算な空間は可分なので第二可算な多様体が可分であることもよくわかる。
\par まず部分空間の開集合をもとの空間の開集合に持ち上げるための次の補題から始まる。

\begin{lem}\label{lem}
空間$X$の部分空間$A$の開集合$U$について
\[
B(U)=\{x\in X\\ |\ d(x,U)<d(x,A\setminus U)\}
\]
と置くと、$B(U)$は$X$の開集合であり、$B(U)\cap A=U$を充たす。ここで$d$とは$X$の距離の一つである。\par
さらに$A$の二つの開集合$U,V$について$U\cap V=\O$ならば$B(U)\cap B(V)=\O$である。
\end{lem}
\begin{proof}
やるとわかる。
\end{proof}


\begin{df}[小さい帰納次元]
空間$X$の小さい帰納次元を帰納的$\ind X$に次のように定義する。\par
$\ind X=-1\iff X=\O$とし、$n-1$まで定義されたとして\par
$\ind X\le n\iff$$X$の任意の点$x$と$x$の任意の近傍$U$に対して$x$の近傍$N$が存在して$\ind (\bd N)\le n-1$かつ$N\subset U$を充たす。\par
と定義する。ちなみに$\ind X=n$とは$\ind X\le n$かつ$n-1<\ind X$である。一般に次元を下から評価することは難しい。また、この次元は位相不変量になっていることはよくわかる。
\end{df}
次の事実\ref{00}は定義からすぐにわかるが、面倒なので証明せずに事実として扱うことにする。
\begin{fact}\label{00}
空間$X$が$0$次元であるとは開かつ閉な集合からなる開基を持つことと同値である。また、$X$は第二加算なのでこの開かつ閉な集合からなる開基は高々可算な族であるとして良い
\end{fact}
次の事実も定義から帰納的に証明することができるはずである。
\begin{fact}\label{fact}
\[
\ind(\rr^n)\le n\ \ind (S^n)\le n
\]
ここで$S^n$は$n$次元球面である。
\end{fact}
\begin{cau}
実は$\ind(\rr^n)=\ind (S^n)=n$が成り立つのであるが、最終的な目標のためには次元が$n$より小さいことが分かっていれば十分である。
\end{cau}
次の事実は可分距離空間の次元論における重要な定理である。
\begin{fact}[分解定理]\label{qq}
空間$X$が$\ind X\le n$とすると$X$の$n+1$個の部分空間$A_0,A_1,\dots ,A_n$が存在して
\[
\ind A_i\le 0
\]
及び
\[
X=\bigcup_{i=0}^nA_i
\]
を充たす。
\end{fact}

\begin{prop}
第二可算な$n$次元位相多様体$M$について
\[
\ind M\le n
\]
\end{prop}
\begin{proof}
局所ユークリッド性から境界が球面と同相になるような開集合が開基を成すことが分かる。この事と事実\ref{fact}からわかる。
\end{proof}

\begin{thm}\label{main}
空間$X$は$\ind X\le n$を充たすとする。そして$\mathcal{U}$を空間$X$の開被覆とする。このとき$\mathcal{U}$の細分である$X$の局所有限な開被覆$\mathcal{V}$と開集合からなる（被覆とは限らない）$n+1$個の族$\mathcal{W}_0,\mathcal{W}_1,\dots,\mathcal{W}_n$が存在して
\[
\mathcal{V}=\mathcal{W}_0\cup\mathcal{W}_1\cup\cdots\cup\mathcal{W}_n
\]
及び
\[
\text{各$i$について$\mathcal{W}_i$は互いに素な族である。}
\]
を充たす。
\end{thm}
\begin{proof}
空間$X$は距離空間なので特にパラコンパクト\footnote{一般に距離空間はパラコンパクトであるが、可分距離空間のパラコンパクト性は一般的なそれよりも容易に証明できる。}であるので$\mathcal{U}$の細分である局所有限な開被覆である$\mathcal{C}$が存在する。ところで$X$は可分距離空間なので特にリンデレフであるから$\mathcal{C}$は可算であるとしても一般性を失わない。$\mathcal{C}=\{U_k\}_{k\in \nn}$としよう。\par
さて事実\ref{qq}から$n+1$個の高々$0$次元部分空間$\{A_i\}_{i=0}^n$が存在して$\dpst X=\bigcup_{i=0}^nA_i$を充たす。\par
ここで各$i$について$A_i$上で$A_i$の被覆$\{A_i\cap U_k\}_{k\in \nn}$を考えよう。$A_i$は$0$次元なので開かつ閉な開基を持つ。また可分でもあるので$\{A_i\cap U_k\}_{k\in \nn}$の細分で高々可算な開かつ閉な集合からなる$A_i$の被覆$\{O_l^{(i)}\}_{l\in \nn}$が存在する。これを用いて
\[
G_1^{(i)}=O_1^{(i)},\ G_j^{(i)}=O_j^{(i)}\setminus \bigcup_{l<j}O_l^{(i)}\ (j\ge 2)
\]
と定義すると$\{G_j^{(i)}\}_{j\in \nn}$は$A_i$の開集合からなる\footnote{$A_i$の閉集合でもあるが、そのことは使わない。}$A_i$の被覆で$\{A_i\cap U_k\}_{k\in \nn}$の細分になる。また定義から明らかに$a\neq b$ならば$G_a^{(i)}\cap G_b^{(i)}=\O$であるので$\{G_j^{(i)}\}_{j\in \nn}$の元は互いに素である。そして$\{I_a^{(i)}\}_{a\in \nn}$を
\[
I_{1}^{(i)}=\{j\in \nn\ |\ G_j\subset U_1\}
\]
とし、$m-1$まで$I_a^{(i)}$が定義されてるとして
\[
I_m^{(i)}=\{j\in\nn\ |\ G_j\subset U_m\}\setminus \bigcup_{a=1}^{m-1}I_{a}^{(i)}
\]
と帰納的に定義しよう。そしてこれを用いて
\[
P_a^{(i)}=\bigcup_{s\in I_a}G_s^{(i)}
\]
と定義しよう。するとこれは$A_i$の開集合であり、$\{P_a^{(i)}\}_{a\in \nn}$は$A_i$の被覆となり、定義から明らかに$\{P_a^{(i)}\}_{a\in \nn}$は互いに素な族である。\par
これを用いて$Q_m^{(i)}=B(P_m^{(i)})\cap U_m$と定義しよう。（補題\ref{lem}参照）このとき補題\ref{lem}から$\{Q_m^{(i)}\}_{m\in \nn}$は互いに素な族となり、$A_i$を被覆する。（ただし$Q_m^{(i)}$は空である場合もある。）また、$Q_m^{(i)}\subset U_m$から$\{U_i\}_{i\in \nn}$の細分になっている。さて、任意の点$x\in X$について$\mathcal{C}$の局所有限性を担保する$x$の近傍を$N$とすると、一般に
\[
N\cap Q_m^{(i)}\neq\O\To N\cap U_m\neq\O
\]
が成り立つことから$\{Q_m^{(i)}\}_{m\in \nn}$は局所有限な族であることが分かる。（ただし$X$の被覆であるとは限らない）\par
各$i=0,1,\dots ,n$ごとに存在する$\{Q_m^{(i)}\}_{m\in \nn}$を$\mathcal{W}_{i}$とすれば$\mathcal{W}_i$は互いに素な族であり、かつ局所有限であり、かつ$\mathcal{C}=\{U_k\}_{k\in \nn}$の細分となっている。そして$\mathcal{V}=\mathcal{W}_0\cup\mathcal{W}_1\cup\cdots\cup\mathcal{W}_n$と置けば各$\mathcal{W}_i$が局所有限で$\mathcal{C}=\{U_k\}_{k\in \nn}$の細分であることから$\mathcal{V}$は局所有限であり、$\mathcal{C}=\{U_k\}_{k\in \nn}$の細分となる。また$A_i\subset \bigcup\mathcal{W}_i$で$\dpst X=\bigcup_{i=0}^nA_i$なので$\mathcal{V}$は$X$の被覆となる。$\mathcal{C}=\{U_k\}_{k\in \nn}$が$\mathcal{U}$の細分であることから$\mathcal{V}$は$\mathcal{U}$の細分であることが分かる。\par
以上で証明は終わった。
\end{proof}











































	
\begin{thebibliography}{9}
\bibitem{-1}yamyamtopさんの次元論PDF,https://yamyamtopo.files.wordpress.com/2015/01/dimensiontheory1.pdf,\ 2016年6月閲覧
\bibitem{0}志賀浩二,多様体論,岩波書店,1976
\bibitem{1}W.Hurewicz and H.Wallman,\textit{Dimension Theory},Princeton Univ. Press,1948
\end{thebibliography}




			\end{document}



\begin{comment}
[可換図式]
\[\xymatrix{
&   & (2) \ar@{=>}[rd]&  &  &  & (5) \ar@{=>}[rd]&\\ 
& (1)\ar@{=>}[ru] \ar@{=>}[rd]&  &(4)\ar@{=>}[r]&(7)& (1)\ar@{=>}[ru] \ar@{=>}[rd]& & (7)&\\ 
&   & (3) \ar@{=>}[ur]&  &  &  &(6) \ar@{=>}[ur]&
}\]

[\bigin \endを使わない方法]

{\prop[aa] aaa}
で
\begin{prop}[aa]
aaa
\end{prop}
と同じ\df \thm \proof でも同様

[直和]
$\amalg$

[同値関係でよく使う波線]
\sim

[引数付きの新命令]
\newcommand{\prn}[1]{\|#1\|_p}

[番号のリセット]

\setcounter{equation}{0}


[字体の変更]

{\rm hogehoge}と使う。

\rm ローマン
\it イタリック
\sl 斜体

[supp]

\newcommand{\supp}{\text{{\rm supp}}}

[中央揃えの改行]

	\begin{eqnarray*}
		hoge1&=&hogehoge
		hoge2&=&hogehofe

	\end{eqnarray*}

&&の部分で揃えられる。


[\tag]
\begin{equation}
f(x) = ax^2 + bx + c \tag{1.2.3}
\end{equation}



[場合分け]

	\begin{cases}
		hoge1 & \text{状況１}\\
		hoge2 & \text{状況２}\\
		・
		・
		・
	\end{cases}



\end{comment}





